\chapter{Visuelle Signale}
\label{chap:signale}

% Inhalt laut Vorgabe:
%  - Dokumentation der visuellen Signale mit Bedeutung

Modul~1 zeigt seinen Zustand über eine Signalsäule mit drei Lampen an
(rot, orange, grün). Die Lampen werden beim Zustandswechsel über das Topic
\texttt{MV/VPS/Delivery/Lamps} gesetzt.

\begin{table}[htbp]
  \centering
  \renewcommand{\arraystretch}{1.5}
  \begin{tabularx}{\textwidth}{@{}l c c c X@{}}
    \toprule
    \textbf{Zustand} & \textbf{Rot} & \textbf{Orange} & \textbf{Grün} & \textbf{Bedeutung für den Bediener} \\
    \midrule
    INIT     & \lampe{lamperot}{1} & \lampe{lampeorange}{1} & \lampe{lampegruen}{1} & Lampentest beim Start \\
    BEREIT   & \lampe{lamperot}{0} & \lampe{lampeorange}{1} & \lampe{lampegruen}{0} & Wartet auf Material im Silo bzw. Freigabe von Modul~2 \\
    VORLAUF  & \lampe{lamperot}{0} & \lampe{lampeorange}{1} & \lampe{lampegruen}{1} & Druckluft läuft vor, Förderband startet gleich \\
    LÄUFT    & \lampe{lamperot}{0} & \lampe{lampeorange}{0} & \lampe{lampegruen}{1} & Normale Förderung \\
    NACHLAUF & \lampe{lamperot}{0} & \lampe{lampeorange}{1} & \lampe{lampegruen}{1} & Förderung endet, Band und Luft laufen aus \\
    STOP     & \lampe{lamperot}{1} & \lampe{lampeorange}{0} & \lampe{lampegruen}{0} & Sicherheitsstopp: Not-Aus (M1 oder M2) oder Trichter voll \\
    \bottomrule
  \end{tabularx}
  \caption[Lampenzustände von Modul~1 und ihre Bedeutung]{Lampenzustände von Modul~1 und ihre Bedeutung
    (\protect\lampe{gray}{1}~an, \protect\lampe{gray}{0}~aus)}
  \label{tab:lampen}
\end{table}

VORLAUF und NACHLAUF zeigen dasselbe Lampenbild, weil in beiden Fällen
Druckluft und gegebenenfalls das Band laufen, aber keine reguläre Förderung
stattfindet.
